\section{Semantyczne reprezentacje tekstu}
W ostatnich latach modele głębokiego uczenia zrewolucjonizowały przetwarzanie
języka naturalnego. W kontekście generowania osadzeń tekstowych,
technologie te umożliwiły znaczący postęp w jakości reprezentacji
semantycznych zdań. Poniżej przedstawiono szczegółowy opis trzech
kluczowych metod omawianych w artykule \cite{DBLP_Gurevych}, z uwzględnieniem ich zalet, ograniczeń oraz zastosowań.

\subsection{BERT (Bidirectional Encoder Representations from Transformers)}
Model \textbf{BERT} stanowi przełom dzięki dwukierunkowemu przetwarzaniu
tekstu, co umożliwia pełne uchwycenie kontekstu zdania. Dzięki warstwom
transformatorowym \textbf{BERT} ustanowił nowe standardy w takich zadaniach jak
klasyfikacja zdań czy analiza podobieństwa tekstowego. Jego
zastosowanie wymaga jednak jednoczesnego podania obu zdań do sieci, co
prowadzi do ogromnych kosztów obliczeniowych.

Dodatkowo \textbf{BERT} nie generuje osadzeń zdań w sposób natywny.
Stosowane podejścia, takie jak uśrednianie danych wyjściowych z warstwy końcowej, często prowadzą do
wyników gorszych od metod tradycyjnych, takich jak \textbf{GloVe}. Zastosowanie
\textbf{BERT} w zadaniach takich jak klasteryzacja czy wyszukiwanie semantyczne
jest ograniczone ze względu na niską jakość generowanych osadzeń oraz
bardzo wysokie wymagania obliczeniowe.

\subsection{RoBERTa (Robustly Optimized BERT Approach)}
\textbf{RoBERTa} to ulepszona wersja \textbf{BERT}, zoptymalizowana pod kątem wydajności w
zadaniach z nadzorem. Dzięki dłuższemu treningowi, większej ilości
danych i lepszej optymalizacji hiperparametrów, \textbf{RoBERTa} osiąga lepsze
wyniki w zadaniach klasyfikacyjnych. Usunięcie tokenów w procesie pretreningu pozwoliło na zwiększenie efektywności
modelu.

Jednak w kontekście generowania osadzeń zdań \textbf{RoBERTa} wykazuje wyniki
porównywalne z \textbf{BERT}.

\subsection{SBERT (Sentence-BERT)}
\textbf{SBERT} unika ograniczenia \textbf{BERT} w generowaniu osadzeń zdań.
Wykorzystuje sieci, które umożliwiają jednoczesne przetwarzanie dwóch zdań za pomocą wag modelu. Osadzenia do wyników są porównywalne, co pozwala na szybkie określenie relacji
semantycznych między zdaniami.

Jednym z kluczowych usprawnień \textbf{SBERT} jest wprowadzenie \textbf{warstwy łączącej}, która uśrednia wartości tokenów, wybiera wartości
maksymalne. Ten model został trenowany na zbiorach NLI (Natural Language Inference) oraz STS.
Dzięki temu \textbf{SBERT} znacząco redukuje czas obliczeń - analiza 10 000 zdań
zajmuje wielokrotnie szybciej w porównaniu do \textbf{BERT}.

\textbf{SBERT} przewyższa jakością osadzenia generowane przez \textbf{BERT}, łącząc
precyzję i elastyczność modelu \textbf{BERT} z wysoką efektywnością obliczeniową.
Dzięki zdolności do szybkiego tworzenia wysokiej jakości osadzeń, \textbf{SBERT}
znajduje zastosowanie w zadaniach takich jak klasteryzacja, wyszukiwanie
semantyczne oraz analiza podobieństwa tekstowego.